% =============================================================================
% Chapter 3: Research Question and Strategy
% Input-ready fragment for the lean toy-model research proposal.
% This file intentionally contains no preamble or document environment.
% =============================================================================

\chapter{Research Question and Strategy}
\label{ch:research-question-strategy}

The ontology in \cref{ch:plain-english-ontology} defines one proposed material
picture. This chapter defines how that picture will be tested. The program is
not organized to accumulate separate analogies for familiar phenomena. It is
organized to determine whether one controlled brane--bulk medium can satisfy the
full set of light, force, particle, and reservoir requirements at the same time.
That distinction controls the order of work, the use of provisional sources,
the treatment of parameters, and the interpretation of both positive and
negative results.

\section{Central research question and scope}

The central research question is:

\begin{quote}
Does there exist at least one mathematically coherent one-medium brane--bulk
system whose common field content, constitutive laws, dynamical branch,
boundary conditions, and internal reservoir structure can simultaneously
support the required far-field sectors, admit stable particle-like throats that
source and respond to those sectors consistently, and close the global material
cycle without independent sector-by-sector retuning?
\end{quote}

This is an existential compatibility test. Let \(\Lambda_S\) denote the region
of admissible common-medium choices that satisfies every mandatory condition
imposed by sector \(S\). The proposed architecture survives the far-field
compatibility stage only if

\[
\boxed{
\Lambda_{\mathrm{common}}
=
\bigcap_{S\in\mathcal S_{\mathrm{mandatory}}}\Lambda_S
\neq\varnothing.
}
\]

The intersection, rather than the union of separately successful regions, is
the research object. A light calculation at one parameter point, an electric
calculation at another, and a gravity calculation on a different dynamical
branch do not together establish one medium. Nor does the absence of an
immediate contradiction establish that the intersection is nonempty. The
program must construct and test at least one common realization.

The scope is deliberately bounded. This is a classical toy-model research
program. It may postulate an arena, material states, constitutive terms, and a
permitted operator basis, then test the consequences of those choices. Even a
fully successful program would establish the existence of a coherent classical
analog only within its declared regimes. It would not by itself establish that
the real universe is literally constituted this way, reproduce exact general
relativity or exact electromagnetism, or supply quantum theory, photon
statistics, or a complete particle spectrum.

Six epistemic distinctions are sufficient for the proposal. A
\emph{postulate} defines the candidate architecture. An \emph{operational
definition} fixes how a brane observer measures or calibrates a quantity. A
\emph{target relation} is an output the model must produce, not an answer that
may be inserted into the equations. A result \emph{derived under stated
assumptions} remains limited to its branch, boundary ensemble, approximation,
and regime. An \emph{open hypothesis} identifies an unresolved mechanism or
solution. A \emph{failure criterion} states what result rejects a named branch,
candidate, mechanism, or architecture. The canonical ontology and closure
ledger control those statuses in detail; this proposal uses them only to keep
the claims proportionate to the evidence.

\section{Capability-first, far-field-first reasoning}

The constitutive structure of the final parent medium is intentionally not
fully frozen at the start. Freezing detailed equations before the required
capabilities are understood would risk spending substantial effort on a stable
slab or nonlinear throat that cannot support a Coulomb carrier, cannot guide
transverse light, or necessarily opens a fatal loss channel. The opposite risk
is equally serious: adding a new term whenever a later sector becomes difficult
would produce several tailored effective theories rather than one medium.

The program therefore begins by deriving \emph{capability contracts}. Each
contract works backward from an observer-facing behavior to the minimum source,
response, symmetry, spectral, conservation, and stability properties that any
successful common medium must possess. It remains broad enough to compare more
than one possible realization, but concrete enough to expose incompatible mode
counts, signs, stiffnesses, characteristic speeds, couplings, reservoir costs,
or boundary conditions.

The far field is the useful first abstraction because it separates long-range
behavior from the unresolved nonlinear core. At distances large compared with
a throat radius and, where relevant, the brane thickness, a localized object can
initially be represented by provisional source data coupled to the background
response operators. This permits calculation of required ranges, source
parities, Green-function structure, retarded poles, force signs, shared
normalizations, leakage, radiation thresholds, and unwanted modes before a
complete throat exists.

The provisional source is a disciplined placeholder, not a solved particle. Its
amplitude, form factor, conservation law, universality, and relation to an actual
throat remain later tests. A far-field result can therefore establish that a
particular source \emph{would require} a healthy odd \(k^2\)-stiffness branch, a
specified transverse pole, or an orientation-even gravity response. It cannot
establish that the common medium admits a stable throat with the assumed source
profile. That gap is closed only after the common background survives and the
nonlinear throat program replaces the placeholders with solved objects.

The strategy has three broad phases. First, light and the far-field force
sectors are converted into compatible requirement statements. Second, those
requirements are intersected and realized, if possible, by one frozen Candidate
Parent System whose stable brane and complete spectrum are solved and from which
every far-field sector is rederived. Third, only a surviving candidate advances
to self-consistent throats, moving and accelerating solutions, two-body contact,
and the local--global reservoir fixed point. Compatibility on paper is not
existence, and existence of a background is not existence of a particle.

\section{Requirement-discovery order and mathematical closure order}

Two orders of work must remain distinct. The \emph{requirement-discovery order}
is the sequence chosen to learn what the common medium must do. The
\emph{mathematical closure order} is the dependency chain that determines what
must actually exist before a later claim is valid for one candidate. Calendar
scheduling may place conditional calculations in parallel, but it does not
alter either logical order.

The proposal uses the following discovery sequence:

\begin{enumerate}[leftmargin=2.2em,itemsep=0.35em]
  \item guided light and the stronger localized photon-packet hypothesis;
  \item gravity as the combined response to localized one-way throat drainage
        and separate distributed porous-brane \(S_{\mathrm{leak}}\) return;
  \item static electricity and oriented charge;
  \item magnetism from a moving oriented source;
  \item accelerating charge and dynamic electromagnetic closure;
  \item passive response, inertia, and low-energy universality;
  \item the complete audit of extra modes, speeds, leakage, radiation, and drag;
  \item cross-sector synthesis and the common viable-region test.
\end{enumerate}

Light leads because it tests the host slab and its shear-supporting spectrum
without requiring a solved particle core. Gravity then establishes the
source-versus-probe distinction and introduces the burden of keeping localized
one-way throat drainage distinct from distributed porous-brane
\(S_{\mathrm{leak}}\) return and global reservoir response.
Static electricity adds the reflection-odd long-range sector; magnetism and
radiation test whether that sector belongs to one conserved dynamical source
hierarchy. Passive and inertial response then test the complete dressed object,
while the full spectral audit collects consequences that may cross several
sectors at once.

The mathematical closure order instead follows the actual dependency chain:

\[
\boxed{
\begin{aligned}
\text{parent dynamical framework}
&\longrightarrow \text{stable finite brane}\\
&\longrightarrow \text{complete background spectrum}\\
&\longrightarrow \text{self-consistent supported throat}\\
&\longrightarrow \text{throat stability and internal spectrum}\\
&\longrightarrow \text{moving and accelerating throat family}\\
&\longrightarrow \text{two-throat interactions}\\
&\longrightarrow \text{global material-cycle and reservoir closure}.
\end{aligned}
}
\]

A later stage cannot confer existence on a missing predecessor. A photon
envelope may be explored on a provisional slab, or a far-field pair kernel may
be studied with provisional sources, but those results remain conditional until
the same profiles and coefficients arise from the frozen background and throat
solutions. The final local and global stages form an iteration rather than a
one-way dependency chain: localized one-way throat drainage changes the
reservoir, while the changed reservoir changes the distributed porous-brane
\(S_{\mathrm{leak}}\) return profile and the asymptotic data under which the
throat must be solved again; it does not return material through the throat.

\section{Compact sector-analysis pattern}

Every sector chapter follows one compact research pattern. It should answer:

\begin{enumerate}[leftmargin=2.2em,label=\textbf{\arabic*.},itemsep=0.42em]
  \item What behavior must a brane observer see, and in what regime?
  \item What material mechanism is proposed to produce that behavior?
  \item What source and response calculation will test the proposal?
  \item What static and dynamic behavior is required from the relevant branch?
  \item What unwanted modes, leakage, radiation, drag, or departures must be
        calculated from the same mechanism?
  \item What does the result require or forbid in the common medium?
  \item What outcome passes, conditionally survives, reopens the work, or
        rejects the branch?
\end{enumerate}

This pattern prevents a sector from passing by reproducing only its desired
leading asymptotic behavior. A static \(1/R\) potential does not by itself close
the dynamics of its carrier. A correct transverse mode count does not establish
finite-slab confinement or a localized photon. A plausible moving-source tensor
does not establish its normalization from the electric sector. Each chapter
must therefore distinguish equilibrium static operators from retarded dynamic
response, and it must include every coupled branch activated by the proposed
source.

The detailed calculation may use a Hessian and an ensemble-appropriate energy,
a cycle-averaged stress for a conservative periodic solution, a retarded
operator for radiation or damping, or a full nonlinear time evolution. The
method is allowed to vary with the physical problem. What remains common is the
handoff: each sector returns a normalized set of required fields, coefficients,
speeds, sources, constraints, exclusions, and failure conditions for the
cross-sector synthesis in \cref{ch:cross-sector-integration}.

\section{Frozen parameters and the proposal--ledger cross-check}

The common-medium claim is meaningful only if shared inputs retain the same
identity and value across dependent calculations. The governing rule is:

\begin{quote}
Once a constitutive term, coefficient, normalization, sign convention,
dynamical branch, source definition, boundary ensemble, or derived upstream
quantity has been fixed for a controlled baseline, it is inherited by every
dependent calculation. A substantive change creates a new version and reopens
all actual downstream dependencies.
\end{quote}

Freezing is not a claim that an input is true, unique, or fundamental. It means
only that the project is testing a named baseline rather than allowing each
sector to select its own preferred theory. During requirement discovery, the
program freezes what a successful medium must do and how the decisive outcome
will be judged. During common-medium realization, Candidate Parent System V1
freezes the actual field inventory, conservative, relaxational, or mixed branch,
equations and constraints, operator basis, symmetries, parameters, boundary and
reservoir data, source conventions, and permitted calibration choices.

A candidate may contain a declared parameter region rather than one numerical
point, but each tested point must be evaluated across every applicable sector.
A favorable result cannot compensate for a ghost, wrong force sign, failed
conservation law, missing mode constraint, nonuniversal coupling, unacceptable
loss channel, or unclosed reservoir. Derived outputs also cease to be adjustable
inputs: a guided-light speed, electric eigenvector, throat source profile, or
reservoir state obtained upstream must be propagated downstream with its stated
uncertainty and regime.

The proposal and derivation ledger provide independent views of the same test:

\begin{center}
\begin{tabularx}{\textwidth}{@{}YY@{}}
\toprule
\textbf{Proposal track} & \textbf{Derivation-ledger track} \\
\midrule
States the observer-facing question and the capability that must be shown
& Starts from declared equations, assumptions, branches, and boundary data \\
Defines the required calculation, unwanted channels, and pass or failure scope
& Performs and records the symbolic, analytic, and numerical calculation \\
Identifies shared quantities and cross-sector dependencies
& Records frozen inputs, outputs, provenance, uncertainty, and reopened descendants \\
Checks whether the mathematics answered the intended physical question
& Checks whether the algebra, implementation, and convergence support the stated claim \\
\bottomrule
\end{tabularx}
\end{center}

Where symbolic verification is feasible, the project uses independently
constructed SymPy and Mathematica calculations to compare invariant results
rather than merely sharing one generated expression. Where symbolic closure is
unavailable, numerical evidence must include branch, boundary, convergence, and
stability checks appropriate to the claim. The detailed artifact and review
procedures remain in the derivation ledger and \cref{ch:verification-milestones-outcomes};
the proposal needs only the principle that no claim is promoted without a
current bottom-up evidence path.

\section{Falsifiability and levels of outcome}

The breadth of the ontology does not protect it from failure. Falsifiability
comes from requiring one frozen system to carry every consequence, separating
static from dynamic questions, calculating unwanted channels, keeping source
properties distinct from probe response, and closing the material and reservoir
ledgers. Words such as ``small,'' ``long lived,'' ``approximately universal,''
or ``acceptable'' must resolve, before decisive evaluation, either to an exact
structural requirement or to a bounded metric with a declared regime and
uncertainty. Mandatory conditions are conjunctive; they are not combined into a
weighted score that permits success in one sector to pay for failure in another.

A negative result must also be attached to the correct level. Failure of one
algebraic implementation or unconverged solver leaves a claim unresolved. Failure
of one approximation rejects that approximation. Failure of a permitted
solution branch rejects that branch. A controlled violation of a mandatory gate
rejects the named Candidate Parent System. The broader architecture fails only
when the mandatory capability intersection is shown to be empty throughout the
admitted theory class, a constitutive mechanism essential to the ontology is
ruled out throughout its permitted realizations, or closure requires prohibited
sector-specific repair.

The principal outcomes are therefore layered:

\begin{center}
\begin{tabularx}{\textwidth}{@{}L{0.27\textwidth}X@{}}
\toprule
\textbf{Outcome} & \textbf{Meaning} \\
\midrule
Sector result
& A named behavior or no-go statement is established within its assumptions. It
  remains valuable without establishing the complete model. \\
Partial analog
& Several sectors survive, or a weaker claim survives after a stronger one
  fails; for example, a valid linear light-wave sector may remain even if no
  self-bound photon packet exists. \\
Candidate failure
& One frozen Candidate Parent System violates a mandatory requirement or cannot
  close its own background, throat, interaction, or reservoir problem. \\
Architecture failure
& The common mandatory region is demonstrably empty across the admitted
  architecture, or success requires untracked new mechanisms or independent
  retuning. \\
Full program survival
& One current parent system passes the common far-field intersection, stable
  background and spectrum, supported-throat, moving, pair, and global-reservoir
  stages within the declared regimes. \\
\bottomrule
\end{tabularx}
\end{center}

An unresolved calculation is neither a pass nor a failure. It should identify
the missing object: a stable slab, source profile, boundary ensemble, dynamic
operator, convergence result, acceptance threshold, or global fixed point.
Likewise, a positive result authorizes only the claim actually tested. Even full
program survival would establish a controlled classical analog architecture,
not empirical identification with nature. Conversely, an empty common region,
an unavoidable extra mode, a failed supported-throat branch, or an unclosed
reservoir cycle is a legitimate and informative result because it identifies
where the proposed one-medium program ceases to be viable.
